\section*{Chapter \ref{ch:b2:amines} --- Amines and Nitrogen Compounds}

\begin{solution}{exo:b2:amines:1}
Primary; secondary; tertiary; quaternary ammonium salt; primary (\omterm{def:b2:huckel:aromatic}{aromatic}).
\end{solution}

\begin{solution}{exo:b2:amines:2}
Ethanamide $<$ aniline $<$ ammonia $<$ methylamine ($\mathrm pK_a$ of the conjugate acids
$-0.6$, 4.60, 9.26, 10.66).
\end{solution}

\begin{solution}{exo:b2:amines:3}
\textit{N}-Methylcyclohexanimine, \ce{C6H10=N-CH3}: amine and ketone in a weakly acidic medium
(pH 4--5), with removal of the water formed (drying agent or azeotropic distillation).
\end{solution}

\begin{solution}{exo:b2:amines:4}
Bromobenzene; benzonitrile; phenol; 4-hydroxyazobenzene \ce{C6H5-N=N-C6H4-OH} (coupling para to the
\ce{O-} group of the phenoxide).
\end{solution}

\begin{solution}{exo:b2:amines:5}
$\mathrm pK_a = 9.80$: at pH 7 the ratio $[\ce{(CH3)3NH+}]/[\ce{(CH3)3N}] = 10^{2.8} \approx 630$: the
protonated form dominates; lemon juice (pH near 2) protonates practically all of it. Shaking an
organic solution of an amine with dilute aqueous acid transfers the amine, as its ammonium ion,
into the water; base frees it again.
\end{solution}

\begin{solution}{exo:b2:amines:6}
1-(Cyclohex-1-en-1-yl)pyrrolidine. After addition and loss of water the nitrogen carries a
positive charge (an iminium ion) and no hydrogen to lose; a proton is lost from the adjacent
carbon instead, forming the \ce{C=C} bond.
\end{solution}

\begin{solution}{exo:b2:amines:7}
Benzaldehyde with propan-2-amine, or propanone with benzylamine, each with \ce{NaBH3CN} at pH
about 6.
\end{solution}

\begin{solution}{exo:b2:amines:8}
Propiophenone (1-phenylpropan-1-one), \ce{C6H5COCH2CH3}: one addition gives the \omterm{def:b2:amines:imine}{imine} anion,
hydrolysed to the ketone.
\end{solution}

\begin{solution}{exo:b2:amines:9}
Potassium phthalimide + 1-bromohexane: \textit{N}-hexylphthalimide; then hydrazine (or hydrolysis)
frees hexan-1-amine. The alkylated phthalimide nitrogen has no lone pair available (it is
delocalised on two carbonyls) and no \ce{N-H} left: it cannot be alkylated a second time.
\end{solution}

\begin{solution}{exo:b2:amines:10}
Bromine water brominates aniline at the three positions ortho and para to the amino
group: 2,4,6-tribromoaniline. \omterm{def:b2:amines:diazonium}{Diazotisation}, then \ce{H3PO2}, replaces the amino group by H: the three
bromines are now 1,3,5 to each other.
\end{solution}

\begin{solution}{exo:b2:amines:11}
4-Nitroaniline + \ce{NaNO2} + \ce{HCl} at 0--\qty{5}{\degreeCelsius}: 4-nitrobenzenediazonium
chloride. Coupling with phenol in weakly basic solution (phenoxide, strongly activated) at the
position para to \ce{O-}: 4-[(4-nitrophenyl)diazenyl]phenol, \ce{O2N-C6H4-N=N-C6H4-OH}, an orange
dye.
\end{solution}

\begin{solution}{exo:b2:amines:12}
The trimethylammonium group is a poor, bulky leaving group and makes the $\beta$ hydrogens on the
\ce{CH3} side more acidic and more accessible; the base removes a hydrogen from the less
substituted carbon: but-1-ene, the less substituted alkene (Hofmann product).
\end{solution}

\begin{solution}{pb:b2:amines:1}
\textbf{1.} $\mathrm{{}^-O_3S{-}C_6H_4{-}NH_3^+}$: the acidic sulfonic group has protonated the amine.
\textbf{2.} The zwitterion is poorly soluble in water; carbonate deprotonates the ammonium group,
giving the soluble sodium sulfanilate.
\textbf{3.} \ce{NaNO2 + HCl -> HNO2 + NaCl}.
\textbf{4.}
\[
  \mathrm{{}^-O_3S{-}C_6H_4{-}NH_2 + HNO_2 + H^+ \to {}^-O_3S{-}C_6H_4{-}N_2^+ + 2\,H_2O} .
\]
\textbf{5.} The diazonium salt decomposes (to the phenol and \ce{N2}) when warm; nitrous acid
decomposes too.
\textbf{6.} $5.00/173.19 = \qty{28.9}{mmol}$: \qty{28.9}{mmol} of \ce{NaNO2}, \qty{1.99}{g}.
\textbf{7.} Excess nitrous acid would react with the coupling partner; it is detected with starch--iodide
paper and destroyed with urea or sulfamic acid.
\textbf{8.} Its ring is strongly activated by the dimethylamino group, a powerful donor.
\textbf{9.} Para to \ce{N(CH3)2}: ortho/para direction, and the ortho positions are hindered by the
methyl groups.
\textbf{10.} $\mathrm{{}^-O_3S{-}C_6H_4{-}N{=}N{-}C_6H_4{-}N(CH_3)_2}$.
\textbf{11.} In strong acid the dimethylamino group is protonated: \ce{-NH(CH3)2+} deactivates the ring
and the coupling stops.
\textbf{12.} The product forms as its protonated, partly soluble acid form; in base the sodium
sulfonate, less soluble in the salty solution, precipitates.
\textbf{13.} Its positive charge is delocalised onto the ring and the terminal nitrogen is a poor
electrophilic centre; only rings activated by \ce{-O-}, \ce{-OH} or \ce{-NR2} react.
\textbf{14.} Yellow-orange in base, red in acid.
\textbf{15.} On a nitrogen of the azo group (the protonated form is stabilised by conjugation with
the dimethylamino group).
\textbf{16.} About $\mathrm pK_a \pm 1$: from 3.1 to 4.4 roughly, around $\mathrm pK_a = 3.47$.
\textbf{17.} Two rings and the azo group form a long conjugated system: the gap between the highest
occupied and the lowest empty $\pi$ levels is small (\cref{prop:b2:huckel:gap}) and the absorption
lies in the visible.
\textbf{18.} Protonation changes the distribution of charge in the $\pi$ system and narrows the gap
further: the absorption moves to longer wavelengths, the colour from yellow to red.
\textbf{19.} \qty{173.19}{g/mol}.
\textbf{20.} \qty{327.33}{g/mol}.
\textbf{21.} \qty{28.87}{mmol}.
\textbf{22.} $0.02887 \times 327.33 = \qty{9.45}{g}$.
\textbf{23.} Decomposition of the diazonium salt, incomplete coupling, solubility of the product in
the mother liquor and the washings.
\textbf{24.} $0.80 \times 9.45 = \boldsymbol{\qty{7.56}{g}}$ of methyl orange.
\end{solution}
